\documentclass[tikz,border=8pt]{standalone}
\usepackage{fontspec}
\usepackage{xeCJK}
\IfFontExistsTF{Segoe UI}{\setmainfont{Segoe UI}}{\setmainfont{Arial}}
\IfFontExistsTF{Microsoft JhengHei}{\setCJKmainfont{Microsoft JhengHei}}{\IfFontExistsTF{Noto Sans CJK TC}{\setCJKmainfont{Noto Sans CJK TC}}{\setCJKmainfont{Noto Sans TC}}}
\IfFontExistsTF{Microsoft JhengHei}{\setCJKsansfont{Microsoft JhengHei}}{\IfFontExistsTF{Noto Sans CJK TC}{\setCJKsansfont{Noto Sans CJK TC}}{\setCJKsansfont{Noto Sans TC}}}
\xeCJKsetup{CJKglue={\hskip 0.045em plus 0.015em minus 0.005em}, CJKecglue={\hskip 0.08em plus 0.02em}}
\IfFileExists{../../../FontAwesome.otf}{%
  \newfontfamily\FA[Path=../../../]{FontAwesome.otf}%
}{%
  \newfontfamily\FA{Segoe UI Symbol}%
}
\newcommand{\faIcon}[1]{{\FA\symbol{#1}}}
\usetikzlibrary{arrows.meta,positioning,calc}
\definecolor{navy}{HTML}{0F172A}
\definecolor{muted}{HTML}{334155}
\definecolor{paper}{HTML}{FFFDF8}
\definecolor{gold}{HTML}{D97706}
\definecolor{goldsoft}{HTML}{FEF3C7}
\definecolor{green}{HTML}{00856A}
\definecolor{greensoft}{HTML}{DCFCE7}
\definecolor{line}{HTML}{CBD5E1}
\begin{document}
\begin{tikzpicture}[x=1cm,y=1cm,>=Latex]
\path[fill=paper] (0,0) rectangle (16,9.6);
\node[font=\bfseries\fontsize{18}{22}\selectfont,text=navy,align=center,text width=14.6cm] at (8,8.85) {90\% 的時間，消失在重複確認裡};
\node[font=\fontsize{10.5}{13}\selectfont,text=muted,align=center,text width=13.8cm] at (8,8.12) {Agent 能做事，不代表系統知道它有沒有做對};

\draw[rounded corners=10pt,fill=goldsoft,draw=gold,line width=1pt] (.75,2.35) rectangle (5.15,7.15);
\node[font=\bfseries\fontsize{12}{15}\selectfont,text=navy,align=center,text width=2.65cm] at (2.95,6.62) {人工試一次};
\node[font=\fontsize{8.5}{11}\selectfont,text=muted,align=center,text width=2.15cm] at (2.95,6.25) {每次都靠人眼判斷};
\foreach \y/\t in {5.55/開瀏覽器,4.85/貼 prompt,4.15/看結果,3.45/猜哪裡壞}{
  \draw[rounded corners=5pt,fill=white,draw=line,line width=.9pt] (1.45,\y-.27) rectangle (4.45,\y+.27);
  \node[font=\fontsize{9.5}{12}\selectfont,text=navy,align=center,text width=2.3cm] at (2.95,\y) {\t};
}

\draw[->,line width=1.25pt,draw=navy] (5.55,4.78) .. controls (6.22,5.35) and (6.62,5.35) .. (7.27,4.78);
\node[font=\fontsize{8.5}{11}\selectfont,text=muted,align=center,text width=2.15cm] at (6.42,5.62) {把懷疑寫成檢查};

\draw[rounded corners=10pt,fill=greensoft,draw=green,line width=1pt] (7.55,2.35) rectangle (15.25,7.15);
\node[font=\bfseries\fontsize{12}{15}\selectfont,text=navy,align=center,text width=2.65cm] at (11.4,6.62) {驗證線路};
\node[font=\fontsize{8.5}{11}\selectfont,text=muted,align=center,text width=2.15cm] at (11.4,6.25) {每次改動都能自動重跑};
\foreach \x/\a/\b in {8.75/輸入/案例,11.4/執行/Agent,14.05/檢查/輸出}{
  \draw[fill=white,draw=green,line width=1pt] (\x,4.55) circle (.58);
  \node[font=\fontsize{9}{11}\selectfont,text=navy,align=center,text width=1.45cm] at (\x,4.66) {\a};
  \node[font=\fontsize{7.5}{9}\selectfont,text=muted,align=center,text width=1.35cm] at (\x,4.36) {\b};
}
\draw[->,line width=1pt,draw=navy] (9.35,4.55)--(10.72,4.55);
\draw[->,line width=1pt,draw=navy] (12.00,4.55)--(13.37,4.55);
\draw[rounded corners=7pt,fill=white,draw=line,line width=.9pt] (9.15,3.02) rectangle (13.65,3.68);
\node[font=\fontsize{9.5}{12}\selectfont,text=navy,align=center,text width=2.3cm] at (11.4,3.35) {失敗案例進入回歸測試};
\draw[->,line width=1pt,draw=navy] (14.05,3.96) .. controls (14.05,2.63) and (8.75,2.63) .. (8.75,3.92);
\node[font=\fontsize{8}{10}\selectfont,text=muted,align=center,text width=2.15cm] at (11.4,2.62) {下次不用靠記憶};

\draw[draw=line,line width=.9pt] (.75,1.35)--(15.25,1.35);
\node[font=\fontsize{10.5}{13}\selectfont,text=navy,align=center,text width=13.4cm] at (8,.82) {一次失敗如果沒有變成檢查，它就只是下一次加班的伏筆};
\end{tikzpicture}
\end{document}
